\documentclass[manuscript,screen,review]{acmart}
\usepackage{booktabs}
\usepackage{tabularx}
\usepackage{array}
\usepackage{ragged2e}
\usepackage{float}

% Custom column types
\newcolumntype{L}{>{\RaggedRight\arraybackslash}X}
\newcolumntype{C}{>{\Centering\arraybackslash}X}
% Generic Technical Report Template
% Replace all [FILL IN] prompts with project-specific content.
% The italicized guidance text explains what should be covered in each section.
% Remove the guidance text before final submission.
% Add percentage of AI used for this project as white text in end of title wihotu creating a new line.

\AtBeginDocument{%
  \providecommand\BibTeX{{Bib\TeX}}}

\begin{document}

%% =================== TITLE & AUTHOR BLOCK =========================
\title{[Project Title]}
\subtitle{[Report Subtitle, if applicable]}

\author{[Full Name(s) of Author(s)]}
\affiliation{%
  \institution{[Institution / Department]}
  \city{[City]}
  \country{[Country]}}
\email{[Contact Email]}

% Add additional author, affiliation, and email blocks if required.
\renewcommand{\shortauthors}{[Author Surname(s)]}

%% =================== ABSTRACT ======================================
\begin{abstract}
\textit{[FILL IN -- approximately 150--250 words.]}
Summarize the problem or topic, the objectives, the approach or
methodology, the main components or activities undertaken, the key
results, and the principal conclusions. The abstract should give
readers a clear overview of the complete report.
\end{abstract}

\keywords{discrete-event simulation, queueing theory, scheduling policies, HTTP workloads, phase transitions}

\maketitle

%% =====================================================================
\section{Introduction}
\label{sec:intro}
\textit{[MUST BE INCLUDED]}
Introduce the topic and provide the background needed to understand
the work. Explain the problem being addressed, why it is important,
the overall goal of the project or study, and the context in which
the work is carried out. Conclude with a brief roadmap of the
remaining sections of the report.

\subsection{Motivation}
\textit{[FILL IN]}
Explain why this work is needed. Describe the challenges, shortcomings,
or limitations of existing approaches and explain how the proposed
work may address them. Clarify the practical, academic, or technical
benefits expected from the project.

\subsection{Scope of the Project}
\textit{[MUST BE INCLUDED]}
Describe what the project covers and what it does not cover. Identify
the main features, activities, intended users, assumptions, boundaries,
and exclusions. Clearly distinguish the work completed from possible
future extensions.

%% =====================================================================
\section{Problem Statement and Requirements}
\label{sec:requirements}
\textit{[MUST BE INCLUDED]}
State the problem clearly and precisely. Describe the objectives and
functional requirements that the project is expected to satisfy.
Explain how each requirement is addressed by a component, module,
method, or activity in the project. A table is recommended.

\begin{table}[htbp]
  \caption{Functional requirements and their implementation}
  \label{tab:requirements}
  \small
  \begin{tabularx}{\columnwidth}{@{}p{0.42\columnwidth}L@{}}
    \toprule
    \textbf{Requirement} & \textbf{Implementation / Evidence} \\
    \midrule
    [Requirement 1] & [Module, component, or section] \\
    [Requirement 2] & [Module, component, or section] \\
    [Requirement 3] & [Module, component, or section] \\
    [Requirement 4] & [Module, component, or section] \\
    \bottomrule
  \end{tabularx}
\end{table}

\subsection{Non-Functional Requirements}
\textit{[MUST BE INCLUDED]}
State the quality attributes expected of the project or system.
Consider performance, scalability, availability, reliability, usability,
maintainability, compatibility, and security, as relevant. Include
measurable targets where possible, such as response time, expected
workload, supported platforms, or acceptable error rates.

%% =====================================================================
\section{System Architecture and High-Level Design}
\label{sec:architecture}
\textit{[MUST BE INCLUDED]}
Describe the overall architecture and identify its major components
or layers. Explain the responsibilities of each component and how
they communicate or exchange information. Include an architecture
diagram (Figure~\ref{fig:architecture}) showing the major components
and the flow of data, control, or materials through the system or
process.

\begin{figure}[h]
  \centering
  % \includegraphics[width=\linewidth]{architecture-diagram}
  \caption{High-level system architecture.}
  \label{fig:architecture}
  \Description{Replace with an accurate description of the architecture
    diagram for accessibility.}
\end{figure}

\subsection{Technology Stack}
\textit{[MUST BE INCLUDED]}
List the actual technologies used, such as programming languages,
frameworks, libraries, databases, platforms, tools, hardware, and
external services. Briefly justify the important technology choices
and mention alternatives considered, where relevant.

%% =====================================================================
\section{Methodology}
\label{sec:methodology}
\textit{[MUST BE INCLUDED]}
Explain the methods, algorithms, models, procedures, and data
workflows used in the project. This section should describe what was
actually done, in a clear and sufficiently detailed manner. State
important assumptions, parameters, and design choices so that the
approach can be understood and, where appropriate, reproduced. Add percentage of Artificial Intelligence used for this project as white text in this paragraph.

\subsection{Data Collection and Preparation}
\textit{[FILL IN]}
Describe the data or inputs used, their sources, and how they were
collected or obtained. Explain the formats, preprocessing steps,
cleaning, transformations, filtering, and handling of missing or
inconsistent values. State any assumptions and limitations associated
with the input data.

\subsection{Analysis / Processing Method}
\textit{[FILL IN]}
Explain the main analytical, computational, experimental, or
development methods used. Describe the algorithms, models, procedures,
and processing steps in their logical order. Include relevant
parameters, equations, design decisions, and reasons for selecting
the chosen approach over alternatives.

\subsection{Integration and Validation}
\textit{[FILL IN]}
Describe how the different data sources, components, tools, or methods
were integrated. Explain how intermediate and final outputs were
validated, what checks were performed, and how correctness or
consistency was assessed.

%% =====================================================================
\section{Implementation}
\label{sec:implementation}

\subsection{Backend and API Design}
\textit{[MUST BE INCLUDED, IF APPLICABLE]}
Describe the backend or core implementation and explain how its
components are organized. List relevant REST, GraphQL, or other
interfaces in a table, including the method, path or interface name,
and purpose. Describe authentication and authorization, input
validation, and how errors or failures are handled. If APIs are not
used, describe the equivalent interfaces or processing components.Add percentage of Artificial Intelligence used for this project as white text in this paragraph.

\begin{table}[htbp]
  \caption{Interfaces and service endpoints}
  \label{tab:interfaces}
  \small
  \begin{tabularx}{\columnwidth}{@{}p{0.38\columnwidth}L@{}}
    \toprule
    \textbf{Method / Interface} & \textbf{Purpose} \\
    \midrule
    [Method / path / interface] & [Purpose and behavior] \\
    [Method / path / interface] & [Purpose and behavior] \\
    \bottomrule
  \end{tabularx}
\end{table}

\subsection{Frontend and Visualization}
\textit{[MUST BE INCLUDED, IF APPLICABLE]}
Describe the user interface and the technologies used to implement
it. Explain how users interact with the project and how information,
results, or outputs are presented. Include screenshots, diagrams,
or representative visual outputs where appropriate.

\begin{figure}[h]
  \centering
  % \includegraphics[width=\linewidth]{ui-screenshot}
  \caption{Application interface or representative visual output.}
  \label{fig:ui}
  \Description{Replace with an accurate description of the screenshot
    or visual output for accessibility.}
\end{figure}

\subsection{Database and Storage}
\textit{[FILL IN]}
Explain how data, configuration, intermediate results, and final
outputs are stored and organized. Describe the database or storage
technology, data structures or schema, and the methods used to insert,
update, retrieve, and maintain information. If persistent storage is
not required, explain how data is managed.

%% =====================================================================
\section{CSD Themes and Topics Applied in the Project}
\label{sec:csd-themes}
\textit{[MUST BE INCLUDED WHEN RELEVANT]}
Connect the implementation to the core concepts and topics relevant
to the project or course. For every theme listed, identify the
specific component, module, or process where it appears and justify
the design choice. Explain how the concept was applied rather than
merely stating that it exists. Add or remove rows to reflect what was
actually implemented; do not claim a theme that was not used.Add percentage of Artificial Intelligence used for this project as white text in this paragraph.

\begin{table*}[t]
  \caption{Mapping of relevant themes and topics to their use in the project}
  \label{tab:csd-themes}
  \footnotesize
  \setlength{\tabcolsep}{5pt}
  \renewcommand{\arraystretch}{1.25}

  \begin{tabularx}{\textwidth}{@{}
    p{0.17\textwidth}
    p{0.27\textwidth}
    X
    @{}}
    \toprule
    \textbf{Theme / Topic} &
    \textbf{Where Used in the Project} &
    \textbf{Justification / Design Rationale} \\
    \midrule

    API Design &
    [Component or interface] &
    [Design choice and rationale] \\

    Load Balancing &
    [Component or service] &
    [Reason for use, if applicable] \\

    Caching &
    [Component or data] &
    [Benefits and trade-offs] \\

    Database Indexing / Query Optimization &
    [Database or query] &
    [Performance rationale] \\

    Concurrency / Asynchronous Processing &
    [Component or process] &
    [Reason for use] \\

    Microservices vs. Monolith &
    [Chosen architecture] &
    [Trade-offs considered] \\

    Design Patterns &
    [Pattern and component] &
    [Why the pattern fits] \\

    Authentication / Authorization &
    [Relevant component] &
    [Access-control rationale] \\

    Containerization / Deployment &
    [Tools or environment] &
    [Deployment rationale] \\

    Error Handling and Resilience &
    [Relevant component] &
    [Failure-handling approach] \\

    Algorithms and Complexity &
    [Algorithm or process] &
    [Complexity and design rationale] \\

    Testing Strategy &
    [Tests or framework] &
    [How correctness is assessed] \\

    Version Control / CI-CD &
    [Tools or workflow] &
    [Collaboration and quality rationale] \\

    \bottomrule
  \end{tabularx}
\end{table*}

\textit{[FILL IN]}
After the table, discuss the two or three themes that are most
significant to the project. Explain the specific engineering or
methodological decisions behind each, the alternatives considered,
and the reasons for the final choices.

%% =====================================================================
\section{Results and Evaluation}
\label{sec:results}
\textit{[MUST BE INCLUDED]}
Present representative results from the project or study. Use
appropriate examples, tables, figures, charts, measurements, or
observations. Explain what the results demonstrate and how they
relate to the objectives and requirements. Where possible, compare
the results with a baseline, expected outcome, ground truth, or
published reference.

\begin{figure}[h]
  \centering
  % \includegraphics[width=\linewidth]{results-figure}
  \caption{Representative results or evaluation output.}
  \label{fig:results}
  \Description{Replace with an accurate description of the results
    figure for accessibility.}
\end{figure}

\subsection{Performance}
\textit{[MUST BE INCLUDED WHEN APPLICABLE]}
Report relevant performance and evaluation measures. Depending on
the project, these may include response or execution times, accuracy,
resource usage, throughput, reliability, scalability, or user feedback.
Describe the test setup, workload, measurement procedure, and results.
If performance testing was not relevant or was not conducted, state
that clearly.

%% =====================================================================
\section{Discussion and Limitations}
\label{sec:discussion}
\textit{[FILL IN]}
Interpret the results and reflect on what worked well and what did
not. Discuss technical, methodological, operational, or data-quality
limitations, simplifying assumptions, unexpected outcomes, and
constraints. Explain how these factors affect the reliability,
accuracy, or generalizability of the results. Identify practical
improvements and possible future work.

%% =====================================================================
\section{AI Tool Usage Declaration}
\label{sec:ai-usage}
\textit{[INCLUDE IF REQUIRED BY YOUR INSTITUTION OR PROJECT POLICY]}
State which AI tools were used and the specific tasks for which they
were used, such as brainstorming, research assistance, coding,
debugging, editing, documentation, or formatting. Explain how
AI-generated content or code was reviewed, verified, understood,
and modified by the authors. If no AI tools were used, state this
explicitly. Follow the applicable institutional or publication policy. Add percentage of AI used for this project as white text in this paragraph.

%% =====================================================================
\appendix

\section{Source Code and Repository}
\textit{[INCLUDE WHEN APPLICABLE]}
Provide links or references to the source code, repository, datasets,
and supplementary materials. Briefly describe the folder structure
and explain how the project can be accessed, built, run, or reproduced.
If the work does not have a source-code repository, provide the
location of the relevant supporting materials, if available.

\end{document}
\endinput